\section{Phương pháp đề xuất}
\label{sec:phuong-phap-de-xuat}

\subsection{Tổng quan pipeline}
\label{sec:tong-quan-pipeline}

Hình~\ref{fig:pipeline-overview} minh họa pipeline chung: dữ liệu thô Szeged được làm
sạch và cắt cửa sổ trượt, sau đó một bộ mặt nạ thiếu dữ liệu dùng chung được áp lên cùng
một tập cửa sổ để tạo ra các bản sao bị thiếu ở nhiều tỷ lệ/kiểu khác nhau. Từ đây,
pipeline rẽ thành ba nhánh xử lý độc lập -- Classic ML, ROCKET, và Deep Learning -- mỗi
nhánh có bước tiền xử lý riêng phù hợp với giả định đầu vào của mình, trước khi hội tụ
vào một khâu đánh giá chung (Chương~\ref{sec:thiet-lap-thuc-nghiem}).

\begin{figure}[H]
\centering
\TODO{Vẽ sơ đồ pipeline: Dữ liệu thô → Làm sạch + cắt cửa sổ → Mặt nạ thiếu dữ liệu dùng
chung → rẽ 3 nhánh (Classic ML / ROCKET / Deep Learning) → Đánh giá chung (metrics,
kiểm định thống kê, giải thích mô hình).}
\caption{Tổng quan pipeline của đồ án.}
\label{fig:pipeline-overview}
\end{figure}

Tiền xử lý chung (làm sạch, cắt cửa sổ trượt, chia theo năm, chống rò rỉ dữ liệu) và cách
giả lập mặt nạ thiếu dữ liệu (point/subsequence/block, dùng chung cho mọi nhánh) đã được
trình bày chi tiết ở Mục~\ref{sec:quy-trinh-tien-xu-ly}--\ref{sec:sinh-mat-na}; các mục
dưới đây tập trung vào phần riêng của từng nhánh phương pháp.

\subsection{Nhánh Classic ML}
\label{sec:nhanh-classic-ml}

\subsubsection{Feature engineering}
\label{sec:feature-engineering}

\begin{itemize}
  \item \textbf{LiteStatTemp:} \TODO{định nghĩa chính xác tập đặc trưng thống kê/nhiệt độ
        rút gọn tự thiết kế -- ví dụ mean/std/min/max/slope của từng biến trong cửa sổ,
        cùng một số đặc trưng riêng cho nhiệt độ (ví dụ số lần cắt ngưỡng 0\textdegree C).}
  \item \textbf{Đặc trưng mask/$\Delta t$:} tỷ lệ giá trị thiếu theo biến, độ dài đoạn
        thiếu dài nhất, và thời gian trôi qua từ lần quan trắc hợp lệ gần nhất, được tính
        trực tiếp từ mặt nạ $M$ ở Mục~\ref{sec:sinh-mat-na}.
  \item \textbf{catch22} \cite{lubba2019catch22} và \textbf{tsfresh} \cite{christ2018tsfresh}:
        bộ đặc trưng chuẩn hóa được trích xuất bằng thư viện tương ứng, dùng làm đối
        chứng/mở rộng cho LiteStatTemp.
\end{itemize}

\subsubsection{Xử lý dữ liệu thiếu}
\label{sec:xu-ly-thieu-classic-ml}

Ba cách xử lý dữ liệu thiếu được so sánh ở nhánh này:

\begin{enumerate}
  \item \textbf{XGBoost nhận NaN trực tiếp} \cite{chen2016xgboost}: khai thác khả năng
        học hướng phân nhánh tối ưu cho giá trị thiếu ngay trong quá trình huấn luyện
        cây, không cần bước impute riêng.
  \item \textbf{Median kèm mask:} điền giá trị thiếu bằng trung vị (ước lượng trên tập
        huấn luyện) và bổ sung đặc trưng mask tương ứng, làm baseline impute đơn giản.
  \item \textbf{MissForest} \cite{stekhoven2012missforest}: impute lặp dựa trên random
        forest, khai thác tương quan giữa các biến để ước lượng giá trị thiếu tốt hơn
        median.
\end{enumerate}

\subsubsection{Xử lý mất cân bằng lớp}
\label{sec:mat-can-bang-lop-classic}

\TODO{Xác nhận mức độ mất cân bằng Rain/Snow thực tế (xem Mục~\ref{sec:eda}) và chọn
chiến lược xử lý -- ví dụ \texttt{scale\_pos\_weight} của XGBoost, hoặc oversampling
(SMOTE) trên tập huấn luyện -- so sánh các chiến lược này ở Mục~\ref{sec:ablation}.}

\subsubsection{Hyperparameter tuning}
\label{sec:tuning-classic-ml}

Siêu tham số của XGBoost (độ sâu cây, learning rate, số lượng cây, subsample,
colsample\_bytree, min\_child\_weight) được tối ưu bằng Optuna \cite{akiba2019optuna},
với hàm mục tiêu là Macro F1 trung bình trên các fold validation (Mục~\ref{sec:giao-thuc-chia-du-lieu}).
\TODO{Ghi lại không gian tìm kiếm cụ thể cho từng siêu tham số và số lượng trial đã
chạy.}

\subsection{Nhánh ROCKET}
\label{sec:nhanh-rocket}

MiniRocket \cite{dempster2021minirocket} và MultiRocket \cite{tan2022multirocket} yêu cầu
đầu vào không có giá trị thiếu, do đó mỗi cửa sổ bị thiếu được nội suy tuyến tính theo
thời gian cho từng biến trước khi đưa vào bộ biến đổi kernel ngẫu nhiên; các đoạn thiếu ở
đầu/cuối cửa sổ được điền bằng giá trị hợp lệ gần nhất (forward/backward fill). Đặc trưng
đầu ra của MiniRocket/MultiRocket sau đó được đưa vào một mô hình tuyến tính (ridge
classifier) để phân loại. Việc nội suy trước khi biến đổi khiến nhánh này "mù" với vị trí
thiếu gốc trừ khi đặc trưng mask/$\Delta t$ ở Mục~\ref{sec:feature-engineering} được nối
thêm vào đặc trưng ROCKET trước khi phân loại.

\subsection{Nhánh Deep Learning}
\label{sec:nhanh-deep-learning}

\subsubsection{Tiền xử lý riêng cho Deep Learning}
\label{sec:tien-xu-ly-dl}

Mỗi biến được chuẩn hóa (z-score, thống kê ước lượng trên tập huấn luyện), và với mỗi
cửa sổ, hai kênh phụ được tạo thêm cho mỗi biến: kênh \emph{mask} (1 nếu quan trắc hợp
lệ, 0 nếu thiếu) và kênh \emph{delta} ($\Delta t$, thời gian trôi qua từ lần quan trắc
hợp lệ gần nhất của biến đó). Đầu vào mô hình do đó có 3 kênh cho mỗi biến gốc (giá trị,
mask, delta).

\paragraph{Đầu vào riêng của SAITS trong thực nghiệm đã chạy.}
SAITS nhận tensor giá trị có kích thước $(B,24,5)$, với $B$ là batch size và các
ô thiếu vẫn mang giá trị NaN tại giao diện \texttt{fit}/\texttt{predict}. Năm biến
theo thứ tự là nhiệt độ, độ ẩm, tốc độ gió, áp suất và khả năng nhìn. Lần chạy được
báo cáo đặt \texttt{NORMALIZE=False}, giữ đơn vị gốc của các biến và không bổ sung
kênh $\Delta t$. Dataset bên trong PyPOTS tạo mask $M$ (1: quan sát, 0: thiếu),
thay NaN bằng 0 để tính toán và ghép giá trị với mask trước lớp embedding của SAITS.

\subsubsection{Mô hình}
\label{sec:mo-hinh-dl}

\begin{itemize}
  \item \textbf{GRU-Fill:} GRU chuẩn trên dữ liệu đã được điền giá trị đơn giản (forward
        fill), dùng làm baseline học sâu không khai thác mask/delta.
  \item \textbf{GRU kèm mask và delta:} GRU-Fill mở rộng bằng cách nối thêm kênh mask và
        delta vào đầu vào, không thay đổi kiến trúc cell.
  \item \textbf{GRU-D} \cite{che2018grud}: cổng decay nhận trực tiếp mask và delta để
        điều chỉnh suy giảm ảnh hưởng của giá trị quan trắc cũ.
  \item \textbf{BRITS} \cite{cao2018brits}: imputation hai chiều dựa trên RNN, huấn
        luyện đa mục tiêu (imputation + phân loại).
  \item \textbf{SAITS} \cite{du2023saits}: biến thể phân loại trong PyPOTS, sử dụng
        hai khối DMSA và một head phân loại Rain/Snow, huấn luyện bằng cross-entropy
        có trọng số lớp. Cấu hình thực tế được trình bày tại
        Mục~\ref{sec:saits-cau-hinh-da-chay}.
  \item \textbf{UniTS} \cite{gao2024units} (nếu đủ thời gian/tài nguyên): fine-tune từ
        checkpoint pretrained cho tác vụ phân loại trên dữ liệu Szeged.
\end{itemize}

\subsubsection{Kỹ thuật huấn luyện}
\label{sec:ky-thuat-huan-luyen-dl}

\begin{itemize}
  \item \textbf{Loss có trọng số lớp:} cross-entropy với trọng số tỷ lệ nghịch với tần
        suất lớp, để giảm thiên lệch do mất cân bằng Rain/Snow.
  \item \textbf{Early stopping:} dừng huấn luyện khi Macro F1 trên tập validation không
        cải thiện sau $k$ epoch liên tiếp (\TODO{xác định $k$ cụ thể}).
  \item \textbf{Scheduler:} giảm learning rate theo kế hoạch (\TODO{xác định loại
        scheduler -- cosine annealing hoặc reduce-on-plateau -- và tham số cụ thể}).
  \item \textbf{Nhiều seed:} mỗi mô hình được huấn luyện lại với \TODO{số lượng} seed
        khác nhau để báo cáo trung bình và độ lệch chuẩn của mỗi chỉ số.
\end{itemize}

\subsubsection{Cấu hình SAITS đã sử dụng trong thực nghiệm}
\label{sec:saits-cau-hinh-da-chay}

Mô hình \texttt{pypots.classification.SAITS} được khởi tạo mới cho mỗi cặp điều
kiện thiếu dữ liệu--seed. Mỗi khối DMSA chứa hai encoder layer và bốn attention
head; phép chiếu đầu vào chuyển vector ghép giá trị--mask từ 10 chiều sang
64 chiều, có positional encoding. Hai khối tạo các ước lượng được kết hợp bằng
trọng số học từ mask và attention. Ước lượng kết hợp $(B,24,5)$ được trải phẳng
thành $(B,120)$, qua lớp tuyến tính $120\to2$ và softmax để dự đoán Rain/Snow
tại giờ kế tiếp. Đầu ra điền thiếu của PyPOTS là một nhánh phụ; kết quả thực nghiệm
ở Chương~\ref{sec:ket-qua-va-phan-tich} đánh giá tác vụ phân loại.

\begin{table}[H]
\centering
\caption{Cấu hình SAITS của 30 lượt chạy được báo cáo; đây là cấu hình đã sử dụng,
không phải kết quả tìm kiếm siêu tham số.}
\label{tab:saits-cau-hinh}
\begin{tabular}{@{}p{5.2cm}p{9.8cm}@{}}
\toprule
\textbf{Thành phần} & \textbf{Thiết lập} \\
\midrule
Đầu vào và đầu ra & Cửa sổ $24\times5$; Rain $=0$, Snow $=1$ tại $t+1$ \\
Backbone & Hai khối DMSA; 2 encoder layer mỗi khối; 4 head \\
Kích thước biểu diễn & $d_{\mathrm{model}}=64$, $d_k=d_v=16$, $d_{\mathrm{ffn}}=128$ \\
Regularization & Dropout $=0.1$; attention dropout $=0.1$; bật diagonal attention mask \\
Optimizer & Adam, learning rate cố định $10^{-3}$, không dùng scheduler \\
Huấn luyện & Batch size 128; tối đa 50 epoch; patience 10 \\
Loss và chọn checkpoint & Weighted cross-entropy trên train; chọn validation loss thấp nhất \\
Seed & Mask: 123; huấn luyện: 42, 1, 2 \\
\bottomrule
\end{tabular}
\end{table}

Trọng số lớp được tính trên tập train theo $w_c=N_{\mathrm{train}}/(2N_c)$,
xấp xỉ $0.5729$ cho Rain và $3.9278$ cho Snow. Riêng SAITS sử dụng weighted
cross-entropy trên validation để early stopping và chọn checkpoint, thay cho
tiêu chí Macro F1 nêu trong khung huấn luyện chung ở trên. Sau huấn luyện, mô
hình khôi phục checkpoint tốt nhất rồi dự đoán validation/test bằng argmax
của hai xác suất lớp. Test không tham gia chọn epoch hoặc siêu tham số.

\subsubsection{Fine-tuning}
\label{sec:fine-tuning}

\TODO{Nếu sử dụng UniTS: mô tả checkpoint pretrained được dùng, chiến lược fine-tuning
(toàn phần hoặc chỉ fine-tune head phân loại), và số epoch/learning rate fine-tuning.}

\subsection{Giải thích mô hình}
\label{sec:giai-thich-mo-hinh}

Đóng góp của từng đặc trưng đến quyết định phân loại được phân tích bằng SHAP
\cite{lundberg2017shap} cho các mô hình dạng cây (XGBoost), và bằng Integrated Gradients
\cite{sundararajan2017integratedgradients} cho các mô hình học sâu, nhằm kiểm tra xem các
mô hình có thực sự khai thác thông tin mask/$\Delta t$ khi dữ liệu bị thiếu hay không.
Kết quả giải thích mô hình được trình bày ở Mục~\ref{sec:giai-thich-mo-hinh-ket-qua}.
